% =============================================================================
% Section 5: Dimensionless Framework
% =============================================================================

\section{Dimensionless framework}
\label{sec:dimensionless_framework}

The central objective of this work is to formulate a closed, physics-based dimensionless framework that predicts the internal flow partitioning, convective heat transfer enhancement, and hydraulic penalties in bypass-controlled single-phase immersion cooling of server electronics. Rather than relying on arbitrary empirical polynomials, the framework is derived from the governing balance laws of mass and momentum across parallel flow paths.

\subsection{Governing causal chain}
\label{sec:causal_chain}

The operational behavior of an open-ratio immersion cooling module is described by the deterministic causal sequence:
\begin{equation}
\mathrm{OR} \;\longrightarrow\; \frac{R_{\mathrm{hyd,bypass}}}{R_{\mathrm{hyd,active}}} \;\longrightarrow\; \Phi_{\mathrm{bypass}} \;\longrightarrow\; \mathrm{Re}_{\mathrm{active}} \;\longrightarrow\; \mathrm{Nu}(\mathrm{Re}_{\mathrm{active}}, \mathrm{Pr}, \mathrm{Gz}) \;\longrightarrow\; R_{\mathrm{th}}, \; \Delta p, \; W_{\mathrm{pump}}
\label{eq:causal_chain}
\end{equation}

\subsection{Parallel hydraulic resistance closure}
\label{sec:hyd_resistance_closure}

Consider a chassis of total height $H_{\mathrm{chassis}}$ and width $W$, containing a heat sink of base height $H_{\mathrm{base}}$ and fin height $H_{\mathrm{fin}}(\mathrm{OR}) = (H_{\mathrm{chassis}} - H_{\mathrm{base}})(1 - \mathrm{OR})$. The total volumetric coolant flow rate $Q_{\mathrm{total}}$ splits into an active heat-sink channel flow $Q_{\mathrm{active}}$ and an overhead bypass clearance flow $Q_{\mathrm{bypass}}$:
\begin{equation}
Q_{\mathrm{total}} = Q_{\mathrm{active}} + Q_{\mathrm{bypass}}
\label{eq:flow_split_sum}
\end{equation}

Because both streams share common plenum inlet and outlet manifolds, the static pressure drop across both parallel branches must be equal:
\begin{equation}
\Delta p_{\mathrm{active}} = \Delta p_{\mathrm{bypass}} = \Delta p_{\mathrm{chassis}}
\label{eq:dp_equality}
\end{equation}

The hydraulic resistance of each stream is defined by $\Delta p = R_{\mathrm{hyd}} \cdot \dot{m}^2 / \rho$. In the laminar developing regime, the Darcy-Weisbach channel friction factor is expressed as $f = C_f / \mathrm{Re}$. The hydraulic resistance of the fin-channel matrix is:
\begin{equation}
R_{\mathrm{hyd,active}} = \frac{f_{\mathrm{ch}} L_{\mathrm{sink}}}{2 D_{h,\mathrm{ch}} A_{\mathrm{ch}}^2} + \frac{K_{\mathrm{loss}}}{2 A_{\mathrm{ch}}^2}
\label{eq:R_active}
\end{equation}
where $A_{\mathrm{ch}} = N_{\mathrm{ch}} \cdot s \cdot H_{\mathrm{fin}}(\mathrm{OR})$ is the total free-flow cross-sectional area through the inter-fin passages, $s$ is the fin spacing, and $D_{h,\mathrm{ch}} = 2 s H_{\mathrm{fin}} / (s + H_{\mathrm{fin}})$ is the fin-channel hydraulic diameter.

Similarly, the hydraulic resistance of the overhead clearance bypass passage is:
\begin{equation}
R_{\mathrm{hyd,bypass}} = \frac{f_{\mathrm{gap}} L_{\mathrm{sink}}}{2 D_{h,\mathrm{gap}} A_{\mathrm{gap}}(\mathrm{OR})^2} + \frac{K_{\mathrm{gap}}}{2 A_{\mathrm{gap}}(\mathrm{OR})^2}
\label{eq:R_bypass}
\end{equation}
where $A_{\mathrm{gap}}(\mathrm{OR}) = W \cdot c = W \cdot (H_{\mathrm{chassis}} - H_{\mathrm{base}}) \cdot \mathrm{OR}$.

Equating pressure drops yields the bypass mass fraction $\Phi_{\mathrm{bypass}} = \dot{m}_{\mathrm{bypass}} / \dot{m}_{\mathrm{total}}$:
\begin{equation}
\Phi_{\mathrm{bypass}} = \frac{1}{1 + \sqrt{\frac{R_{\mathrm{hyd,bypass}}}{R_{\mathrm{hyd,active}}}}} = \frac{1}{1 + C_1 \cdot \left(\frac{1 - \mathrm{OR}}{\mathrm{OR} + \epsilon}\right)^m \cdot \mathrm{Re}_{\mathrm{in}}^n \cdot \mathrm{Pr}^k}
\label{eq:phi_bypass_closure}
\end{equation}
where $\epsilon = 10^{-4}$ ensures numerical boundedness as $\mathrm{OR} \to 0$.

\subsubsection{Limiting physical behavior}
Equation~(\ref{eq:phi_bypass_closure}) strictly satisfies the physical boundary constraints:
\begin{align}
\lim_{\mathrm{OR} \to 0} \Phi_{\mathrm{bypass}}(\mathrm{OR}) &= 0 \quad \text{(Fully shrouded, zero tip-clearance leakage)}, \label{eq:limit_or0} \\
\lim_{\mathrm{OR} \to 1} \Phi_{\mathrm{bypass}}(\mathrm{OR}) &= 1 \quad \text{(Fins fully recessed, complete fluid bypass)}. \label{eq:limit_or1}
\end{align}

\subsection{Effective active Reynolds number}
\label{sec:re_active}

The Reynolds number governing convective transport within the active fin channels is not set by the pump flow rate alone, but by the partitioned active mass flux:
\begin{equation}
\mathrm{Re}_{\mathrm{active}} = (1 - \Phi_{\mathrm{bypass}}) \cdot \mathrm{Re}_{\mathrm{in}} \cdot \left(\frac{A_{\mathrm{inlet}}}{A_{\mathrm{ch}}}\right) \cdot \left(\frac{D_{h,\mathrm{ch}}}{D_{h,\mathrm{inlet}}}\right)
\label{eq:re_active_def}
\end{equation}

\subsection{Thermally developing Nusselt number closure}
\label{sec:nu_closure}

Because the dielectric coolants possess high Prandtl numbers ($\mathrm{Pr} = 9.2\text{--}210$), the thermal entrance length $L_{\mathrm{th}} \approx 0.05 \cdot \mathrm{Re}_{\mathrm{active}} \cdot \mathrm{Pr} \cdot D_{h,\mathrm{ch}}$ substantially exceeds the physical length of the heat sink ($L_{\mathrm{sink}} = 80\text{ mm}$). Flow within the fin channels is thermally developing throughout the domain.

The Nusselt number is modeled using an asymptotic composite Graetz formulation:
\begin{equation}
\mathrm{Nu} = \left[ \mathrm{Nu}_{\mathrm{fd}}^3 + \left( C_2 \cdot \mathrm{Gz}^{p} \cdot \left(\frac{\mu_b}{\mu_w}\right)^{0.14} \right)^3 \right]^{1/3}
\label{eq:nu_composite}
\end{equation}
where $\mathrm{Gz} = \mathrm{Re}_{\mathrm{active}} \cdot \mathrm{Pr} \cdot (D_{h,\mathrm{ch}} / L_{\mathrm{sink}})$ is the channel Graetz number, $\mathrm{Nu}_{\mathrm{fd}} \approx 3.66\text{--}4.36$ represents the fully developed asymptotic limit, and $(\mu_b / \mu_w)^{0.14}$ accounts for temperature-dependent property variations across the thermal boundary layer.

\subsection{Overall thermal resistance and pumping power}
\label{sec:rth_and_fom}

The convective heat transfer coefficient is obtained directly from $h = \mathrm{Nu} \cdot k_f / D_{h,\mathrm{ch}}$. The total junction-to-inlet thermal resistance is given by the series sum:
\begin{equation}
R_{\mathrm{th}} = \frac{T_{\mathrm{chip,max}} - T_{\mathrm{inlet}}}{P_{\mathrm{TDP}}} = R_{\mathrm{TIM}} + R_{\mathrm{spread}} + \frac{1}{\eta_o h A_{\mathrm{wetted}}}
\label{eq:rth_sum}
\end{equation}
where $\eta_o = 1 - (N_{\mathrm{fin}} A_{\mathrm{fin}} / A_{\mathrm{wetted}})(1 - \eta_{\mathrm{fin}})$ is the overall surface efficiency, with fin efficiency $\eta_{\mathrm{fin}} = \tanh(m_f H_{\mathrm{fin}}) / (m_f H_{\mathrm{fin}})$ and fin parameter $m_f = \sqrt{2 h / (k_{\mathrm{copper}} t_{\mathrm{fin}})}$.

The pumping power required to drive the flow through the 1U chassis is:
\begin{equation}
W_{\mathrm{pump}} = Q_{\mathrm{total}} \cdot \Delta p_{\mathrm{chassis}}
\label{eq:wpump}
\end{equation}
leading to the global thermal-hydraulic Figure of Merit (FOM):
\begin{equation}
\mathrm{FOM} = \frac{1}{R_{\mathrm{th}} \cdot W_{\mathrm{pump}}} \quad [\mathrm{K}^{-1} \cdot \mathrm{W}^{-1}]
\label{eq:fom}
\end{equation}
